\documentclass[11pt]{scrartcl}
\usepackage{evan}
\usepackage[print]{mathteam}

\title{Notes}

\begin{document}

\maketitle

\begin{abstract}
    It is a maxim of math competitions that memorization matters less than problem-solving skill. However, there are some cases where memorization is important. Here are some things I should remember, but have trouble with.
\end{abstract}

\section{Complex bashing}

\begin{itemize}
    \ii Let the unit circle $\omega$ be tangent to $\ol{BC}$, $\ol{CA}$, and $\ol{AB}$ at $D$, $E$, and $F$, so $\omega$ is an incircle or excircle of $\triangle ABC$. Then the center of $(ABC)$ is $\frac{2(d + e + f)def}{(d + e)(e + f)(f + d)}$, and the intersections between $(DEF)$ and $(ABC)$ are the roots of the two equivalent quadratics
    \begin{align*}
        2m_{11}x^2 - (6m_3 + m_{21})x + 2m_{211} &= 0 \\
        \ol a(2d + a)x^2 - (d^2\ol a + 2d + a + 2da\ol a)x + da(d\ol a + 2) &= 0.
    \end{align*}
    \ii Jacobi bialternant formula: $\frac{a_{\lambda + \rho}}{a_\rho} = s_\lambda$.
    \ii Let $k$ be a field of rational functions of some variables with coefficients in $\QQ$ (for example, $\QQ(a,b,c)$). Let $h \in k[x]$ have roots $\alpha$ and $\beta$. Then for a rational function $\frac fg \in k(x)$ with $\max(\deg f,\deg g) + 1 = n$, we have $\frac{f(\alpha)}{g(\alpha)} = \frac{f(\beta)}{g(\beta)}$ if the $n \times n$ determinant with one row as the coefficients of $f$, one row as the coefficients of $g$, and $n - 2$ rows as the coefficients of the coefficients of $x^ih(x)$ for each $0 \le i < n - 2$ is zero. In particular, if $h$ is irreducible, this is an if and only if. The proof is by interpreting $k[x]$ as a vector space.
    \ii In $\triangle ABC$, the $A$-Dumpty point $D$ is the center of the pair of spiral similarities sending $\ol{AB}$ to $\ol{CA}$ and $\ol{AC}$ to $\ol{BA}$. In terms of complex numbers, $D = \frac{A^2 - BC}{2A - B - C}$. In particular, if $A^2 = BC$, then $D = 0$. It may be checked (by setting $(ABC)$ to be the unit circle) that $D \in (BOC)$, where $O$ is the circumcenter of $\triangle ABC$.
    \ii Let $ABCD$ be a convex quadrilateral inscribed in the unit circle. It is possible to choose complex numbers $a$, $b$, $c$, and $d$ such that $A = a^2$, $B = b^2$, $C = c^2$, $D = d^2$, and the arc midpoints of $\wideparen{AB}$, $\wideparen{AC}$, $\wideparen{CD}$, $\wideparen{DA}$, $\wideparen{ABC}$, and $\wideparen{BCD}$ are $-ab$, $-bc$, $-cd$, $+da$, $+ac$, and $+bd$, respectively.
\end{itemize}

\section{Other}

\begin{itemize}
    \ii For an odd prime $p$ and $0 \le n \le \frac{p - 1}2$, we can transform $\binom{2n}n \equiv (-4)^n\binom{\frac{p - 1}2}n \pmod p$.
    \ii Factorials modulo $p$ are almost always dealt with considering inverses and exploiting the symmetry $x \mapsto -x$. Consider $(\frac{p - 1}2)!$ often.
\end{itemize}

\section{Estimation}

We begin with a log table. Here $\log$ is the natural log.
\begin{center}
\begin{tabular}{c|c}
$\log$ & value \\ \hline
$\log 2$ & $0.693147$ \\
$\log 3$ & $1.09861$ \\
$\log 5$ & $1.60944$ \\
$\log 7$ & $1.94591$
\end{tabular}
\end{center}
You will usually only need the first three significant figures, but sometimes problems are weird.

We continue with some facts about space. I don't believe in significant figures.
\begin{itemize}
    \ii The speed of light is $3 \cdot 10^8$ meters per second.
    \ii One AU is $8$ light-minutes. It takes light $1.3$ seconds to reach the moon.
    \ii The radius of the Earth is $6000$ kilometers.
    \ii Mercury, Venus, Earth, Mars, the asteroid belt, Jupiter, Saturn, Uranus, Neptune, and Pluto are $\frac13$, $\frac23$, $1$, $1.5$, $3$, $5$, $10$, $20$, $30$, and $40$ AU from the Sun, respectively.
    \ii The Sun has a mass of $2 \cdot 10^{30}$ kilograms.
    \ii Jupiter weighs $\frac1{1000}$ as much as the Sun. Saturn weighs $\frac13$ as much as Jupiter. Uranus and Neptune weigh $\frac16$ as much as Saturn, or $10^{26}$ kilograms.
    \ii Earth weighs $\frac1{1000}$ as much as Saturn. Venus weighs as much as Earth. Mars weighs $\frac1{10}$ as much as Earth. Mercury weighs $\frac12$ as much as Mars. The moon weighs $\frac1{100}$ as much as Earth.
    \ii The sun and the gas giants have a density of around $1.4$ grams per cubic centimeter, except for Saturn, which is around half that. The rocky planets have a density of around $5$ grams per cubic centimeter.
    \ii The Milky Way has a diameter of $10^5$ light years. Andromeda is $25$ Milky Way diameters away. If we were to flatten the Milky Way into a uniform disk, we would have $50$ solar masses per square light year.
\end{itemize}

\end{document}